%% GENERATED by generate_papers.py -- edit the bank_*.py files, not this file.
\documentclass[11pt,a4paper]{article}
\usepackage[a4paper, top=20mm, bottom=20mm, left=18mm, right=18mm, headheight=14pt]{geometry}
\usepackage[T1]{fontenc}
\usepackage{lmodern}
\usepackage{amsmath, amssymb, mathtools}
\usepackage{braket}
\usepackage{array, tabularx, booktabs, multicol}
\usepackage{enumitem}
\usepackage[dvipsnames]{xcolor}
\usepackage[most]{tcolorbox}
\usepackage{needspace}
\usepackage{fancyhdr, lastpage}
\usepackage{pifont}
\usepackage{tikz}
\usepackage[hidelinks]{hyperref}
\newcolumntype{L}{>{\raggedright\arraybackslash}X}
\definecolor{brand}{HTML}{1F4E79}
\definecolor{okgreen}{HTML}{1B7F3B}
\setlength{\parindent}{0pt}
\setlength{\parskip}{3pt}
\newcommand{\optlabel}[1]{\textbf{(#1)}}
\newcommand{\correctlabel}[1]{\textcolor{okgreen}{\textbf{(#1)}\,\ding{51}}}
\newenvironment{qitem}[1]{\begin{list}{}{\setlength{\leftmargin}{3.3em}\setlength{\labelwidth}{3.1em}%
  \setlength{\labelsep}{0.2em}\setlength{\topsep}{3pt}\setlength{\parsep}{3pt}}\item[\textbf{Q.#1}\hfill]}{\end{list}}
\newcommand{\qsep}{\par\vspace{2pt}{\color{black!15}\hrule}\vspace{4pt}}
\newcommand{\qmeta}[2]{{\footnotesize\color{black!60}\textsc{#1} \;$\cdot$\; #2}\par}
\newtcolorbox{passagebox}[1]{enhanced, breakable, colback=brand!4, colframe=brand!70, fonttitle=\bfseries,
  title={#1}, left=6pt, right=6pt, top=4pt, bottom=4pt}
\newtcolorbox{sectionbanner}{enhanced, colback=brand, colframe=brand, coltext=white, fontupper=\large\bfseries,
  left=8pt, right=8pt, top=6pt, bottom=6pt, sharp corners}
\newtcolorbox{infobox}[1]{enhanced, breakable, colback=black!3, colframe=black!45, fonttitle=\bfseries,
  title={#1}, left=6pt, right=6pt, top=4pt, bottom=4pt}
\newtcolorbox{solbox}[1]{enhanced, breakable, colback=okgreen!3, colframe=okgreen!60!black, fonttitle=\bfseries,
  title={Answer: (#1) \quad Detailed solution}, left=6pt, right=6pt, top=3pt, bottom=4pt, before skip=4pt, after skip=6pt}
\newcommand{\bub}[1]{\tikz[baseline=-0.65ex]\node[draw, circle, inner sep=0pt, minimum size=4.6mm, font=\scriptsize]{#1};}
\allowdisplaybreaks
\setlength{\emergencystretch}{3em}

\begin{document}
\pagestyle{fancy}\fancyhf{}
\fancyhead[L]{\small\color{black!70}Shor's Algorithm --- Assessment}
\fancyhead[R]{\small\color{black!70}\textbf{Question Paper}}
\fancyfoot[C]{\small Page \thepage\ of \pageref*{LastPage}}
\renewcommand{\headrulewidth}{0.4pt}

\thispagestyle{empty}
\begin{center}
{\color{brand}\rule{\linewidth}{1.2pt}}\\[6pt]
{\Huge\bfseries\color{brand} Shor's Algorithm}\\[4pt]
{\Large Number theory, order finding, phase estimation, implementation and cryptography}\\[4pt]
{\large\bfseries QUESTION PAPER}\\[2pt]
{\color{brand}\rule{\linewidth}{1.2pt}}\\[1pt]
\end{center}
\vspace{-2pt}
\begin{tabularx}{\linewidth}{@{}lXlr@{}}
\textbf{Time allowed:} & 3 hours & \textbf{Maximum marks:} & 240\\
\textbf{Questions:} & 80 (4 sections $\times$ 20) & \textbf{Answer key:} & separate Answer Book
\end{tabularx}
\vspace{6pt}

\textbf{Name of the candidate:} \hrulefill\quad \textbf{Roll No.:} \rule{35mm}{0.4pt}

\vspace{6pt}
\begin{infobox}{General instructions}
\begin{enumerate}[leftmargin=1.4em, itemsep=1pt, topsep=1pt]
\item This question paper has \textbf{four sections} of \textbf{20 questions} each:
\textbf{A}~Mathematical and Number Theory Foundations, \textbf{B}~Quantum Computing Foundations, \textbf{C}~Shor's Algorithm Fundamentals, \textbf{D}~Advanced Shor's Algorithm and Implementation.
\item Every question is of the \textbf{single-correct} type: four options, exactly one correct. Question formats used:
single-correct MCQ, matching lists (List-I/List-II), Assertion--Reason, Statement~I/II and paragraph-based questions.
\item Mark your answers on the \textbf{Response Sheet} at the end of this paper by darkening ONE bubble per question.
\item Calculators are \textbf{not} allowed. Rough work may be done on separate sheets.
\item Unless a question says otherwise, use the \textbf{conventions} printed below; they apply to every section.
\item The detailed solutions are in the separate \emph{Answer Book}; do not consult it before finishing the test.
\end{enumerate}
\end{infobox}
\begin{infobox}{Marking scheme (every question)}
Each question has \textbf{four} options (A), (B), (C) and (D). \textbf{ONLY ONE} of these four options is correct.\\[2pt]
\begin{tabular}{@{}ll@{}}
\emph{Full marks}: & $+3$ if ONLY the correct option is chosen;\\
\emph{Zero marks}: & $\phantom{+}0$ if none of the options is chosen (question unanswered);\\
\emph{Negative marks}: & $-1$ in all other cases (wrong option, or more than one option chosen).
\end{tabular}
\end{infobox}

\begin{infobox}{Conventions and notation used throughout the paper}
\footnotesize
\begin{itemize}[leftmargin=1.2em, itemsep=1pt, topsep=2pt]
\item $a \bmod N$ is the remainder in $\{0, \dots, N - 1\}$; the \textbf{order} $r = \operatorname{ord}_N(a)$ is the least $r > 0$ with $a^r \equiv 1 \pmod N$ (defined for $\gcd(a, N) = 1$);
      $\varphi$ is Euler's totient, $\lambda$ the Carmichael function. ``Factors'' means non-trivial factors ($1 < p < N$).
\item \textbf{Reduction:} if $r$ is even and $a^{r/2} \not\equiv -1 \pmod N$, then $\gcd(a^{r/2} - 1, N)$ and $\gcd(a^{r/2} + 1, N)$ are non-trivial factors; otherwise the attempt fails.
\item \textbf{Circuit:} $t$ counting qubits (Hadamards), work register in $\ket1$, controlled $U_{a^{2^j} \bmod N}$ with $U_c\ket y = \ket{cy \bmod N}$ for $y < N$, inverse QFT, measurement of the
      counting register giving $y \in \{0, \dots, 2^t - 1\}$ (most significant bit first). \textbf{Post-processing:} \texttt{Fraction(y, 2**t).limit\_denominator(N)} gives $s/r'$;
      the order is $r'$ or its smallest multiple $kr'$ with $a^{kr'} \equiv 1$. $y = 0$ gives no candidate.
\item \textbf{Notebook} (\texttt{Shor.ipynb}, defaults): $N = 15$, $a = 7$, $n = $ \texttt{N.bit\_length()} work qubits, $t = 2n$ counting qubits, multipliers as exact
      permutation matrices (\texttt{cudaq.register\_operation}), $1000$ shots. \textbf{Previously recorded run}: CUDA-Q 0.16.0, CPU target \texttt{qpp-cpu}, no GPU; $12$ qubits,
      $51$ gates (plus $8$ measurements), outcomes $0, 64, 128, 192$; $0.766$ of the shots gave $15 = 3 \times 5$. No physical quantum processor was used.
\item State-vector memory: $16$ bytes per amplitude ($2^q$ amplitudes for $q$ qubits). $1$ MB $= 10^6$ bytes, $1$ GB $= 10^9$ bytes.
\end{itemize}
\end{infobox}
\newpage

\begin{sectionbanner}SECTION A \enspace$\cdot$\enspace Mathematical and Number Theory Foundations \hfill Q.1--20\end{sectionbanner}

{\small \textbf{Syllabus:} modular arithmetic, gcd and the Euclidean algorithm, modular inverses, Euler's totient, multiplicative order, repeated squaring, the Chinese remainder theorem, square roots of unity, classical factoring methods and the reduction to order finding.\\ \textbf{Formats:} Single-correct MCQ (Q.1--14); Matching lists (Q.15--16); Assertion--Reason (Q.17); Statement I / II (Q.18); Paragraph based (Q.19--20). Each question: $+3$ / $0$ / $-1$.}
\par\medskip

\needspace{7\baselineskip}
\begin{qitem}{1}
Using the Euclidean algorithm, $\gcd(1071, 462)$ equals
\par\smallskip\noindent\begin{tabularx}{\linewidth}{@{}*{4}{L}@{}}
\optlabel{A}\,$42$ & \optlabel{B}\,$7$ & \optlabel{C}\,$21$ & \optlabel{D}\,$147$
\end{tabularx}
\end{qitem}
\qsep

\needspace{7\baselineskip}
\begin{qitem}{2}
The multiplicative inverse of $3$ modulo $20$ (the $b$ with $3b \equiv 1 \pmod{20}$, $0 < b < 20$) is
\par\smallskip\noindent\begin{tabularx}{\linewidth}{@{}*{4}{L}@{}}
\optlabel{A}\,$3$ & \optlabel{B}\,$13$ & \optlabel{C}\,$17$ & \optlabel{D}\,$7$
\end{tabularx}
\end{qitem}
\qsep

\needspace{7\baselineskip}
\begin{qitem}{3}
Euler's totient $\varphi(91)$ equals
\par\smallskip\noindent\begin{tabularx}{\linewidth}{@{}*{4}{L}@{}}
\optlabel{A}\,$84$ & \optlabel{B}\,$78$ & \optlabel{C}\,$72$ & \optlabel{D}\,$90$
\end{tabularx}
\end{qitem}
\qsep

\needspace{7\baselineskip}
\begin{qitem}{4}
The multiplicative order of $2$ modulo $11$ is
\par\smallskip\noindent\begin{tabularx}{\linewidth}{@{}*{4}{L}@{}}
\optlabel{A}\,$11$ & \optlabel{B}\,$2$ & \optlabel{C}\,$10$ & \optlabel{D}\,$5$
\end{tabularx}
\end{qitem}
\qsep

\needspace{7\baselineskip}
\begin{qitem}{5}
Using repeated squaring, $3^{13} \bmod 17$ equals
\par\smallskip\noindent\begin{tabularx}{\linewidth}{@{}*{4}{L}@{}}
\optlabel{A}\,$5$ & \optlabel{B}\,$3$ & \optlabel{C}\,$16$ & \optlabel{D}\,$12$
\end{tabularx}
\end{qitem}
\qsep

\needspace{7\baselineskip}
\begin{qitem}{6}
Square-and-multiply computes $a^{1000} \bmod N$ from the binary expansion $1000 = 1111101000_2$. The number of modular multiplications
(squarings plus extra multiplications) is
\par\smallskip\noindent\begin{tabularx}{\linewidth}{@{}*{4}{L}@{}}
\optlabel{A}\,$10$ & \optlabel{B}\,$999$ & \optlabel{C}\,$16$ & \optlabel{D}\,$14$
\end{tabularx}
\end{qitem}
\qsep

\needspace{7\baselineskip}
\begin{qitem}{7}
For $N = 35$, which of the following can \emph{not} be the multiplicative order of any $a$ with $\gcd(a, 35) = 1$?
\par\smallskip\noindent\begin{tabularx}{\linewidth}{@{}*{4}{L}@{}}
\optlabel{A}\,$6$ & \optlabel{B}\,$5$ & \optlabel{C}\,$4$ & \optlabel{D}\,$12$
\end{tabularx}
\end{qitem}
\qsep

\needspace{7\baselineskip}
\begin{qitem}{8}
For $N = 15$, a base $a$ is drawn uniformly from $\{2, 3, \dots, 13\}$. The probability that $\gcd(a, 15) > 1$, so that a factor is found classically
without any quantum step, is
\par\smallskip\noindent\begin{tabularx}{\linewidth}{@{}*{4}{L}@{}}
\optlabel{A}\,$7/15$ & \optlabel{B}\,$1/2$ & \optlabel{C}\,$8/15$ & \optlabel{D}\,$1/3$
\end{tabularx}
\end{qitem}
\qsep

\needspace{7\baselineskip}
\begin{qitem}{9}
The unique $x \in \{0, \dots, 14\}$ with $x \equiv 2 \pmod 3$ and $x \equiv 3 \pmod 5$ is
\par\smallskip\noindent\begin{tabularx}{\linewidth}{@{}*{4}{L}@{}}
\optlabel{A}\,$13$ & \optlabel{B}\,$2$ & \optlabel{C}\,$11$ & \optlabel{D}\,$8$
\end{tabularx}
\end{qitem}
\qsep

\needspace{7\baselineskip}
\begin{qitem}{10}
The complete set of solutions of $x^2 \equiv 1 \pmod{15}$ with $0 \le x < 15$ is
\par\smallskip\noindent\begin{tabularx}{\linewidth}{@{}*{4}{L}@{}}
\optlabel{A}\,$\{1, 4, 14\}$ & \optlabel{B}\,$\{1, 14\}$ & \optlabel{C}\,$\{1, 11\}$ & \optlabel{D}\,$\{1, 4, 11, 14\}$
\end{tabularx}
\end{qitem}
\qsep

\needspace{7\baselineskip}
\begin{qitem}{11}
For $N = 21$ and $a = 5$, the order $r$ and the outcome of Shor's classical reduction are
\begin{itemize}[leftmargin=2.6em, labelsep=0.4em, itemsep=2pt, topsep=3pt]
\item[\optlabel{A}] $r = 12$; it gives the factors $3$ and $7$
\item[\optlabel{B}] $r = 6$; the reduction fails because $5^3 \equiv -1 \pmod{21}$
\item[\optlabel{C}] $r = 3$; it fails because $r$ is odd
\item[\optlabel{D}] $r = 6$; it gives the factors $3$ and $7$
\end{itemize}
\end{qitem}
\qsep

\needspace{7\baselineskip}
\begin{qitem}{12}
$2^{100} \bmod 13$ equals
\par\smallskip\noindent\begin{tabularx}{\linewidth}{@{}*{4}{L}@{}}
\optlabel{A}\,$12$ & \optlabel{B}\,$3$ & \optlabel{C}\,$1$ & \optlabel{D}\,$9$
\end{tabularx}
\end{qitem}
\qsep

\needspace{7\baselineskip}
\begin{qitem}{13}
Trial division factors $N = 10403$ by testing divisors $d \ge 2$ in increasing order. The largest $d$ that may have to be tested is
\par\smallskip\noindent\begin{tabularx}{\linewidth}{@{}*{4}{L}@{}}
\optlabel{A}\,$103$ & \optlabel{B}\,$5201$ & \optlabel{C}\,$101$ & \optlabel{D}\,$102$
\end{tabularx}
\end{qitem}
\qsep

\needspace{7\baselineskip}
\begin{qitem}{14}
For an odd $N$ with $k$ distinct prime factors (not a prime power) and $a$ uniformly random among the residues coprime to $N$, the probability
that $r$ is even and $a^{r/2} \not\equiv -1 \pmod N$ is at least $1 - 2^{-(k-1)}$. For $N = 105 = 3\cdot5\cdot7$ this lower bound is
\par\smallskip\noindent\begin{tabularx}{\linewidth}{@{}*{4}{L}@{}}
\optlabel{A}\,$3/4$ & \optlabel{B}\,$1/8$ & \optlabel{C}\,$1/2$ & \optlabel{D}\,$7/8$
\end{tabularx}
\end{qitem}
\qsep

\needspace{7\baselineskip}
\begin{qitem}{15}
Match each quantity for $N = 15$ (List-I) with its value (List-II).
\par\smallskip\noindent{\renewcommand{\arraystretch}{1.35}%
\begin{tabularx}{\linewidth}{@{}c>{\raggedright\arraybackslash}p{0.33\linewidth}@{\hspace{1.2em}}cL@{}}
\toprule
\multicolumn{2}{@{}l}{\textbf{List-I}} & \multicolumn{2}{l@{}}{\textbf{List-II}} \\ \midrule
\textbf{(P)} & $\varphi(15)$ & \textbf{(1)} & $13$ \\
\textbf{(Q)} & order of $4$ modulo $15$ & \textbf{(2)} & $8$ \\
\textbf{(R)} & Carmichael value $\lambda(15) = \operatorname{lcm}(\lambda(3), \lambda(5))$ & \textbf{(3)} & $2$ \\
\textbf{(S)} & $7^{-1} \bmod 15$ & \textbf{(4)} & $4$ \\
\bottomrule
\end{tabularx}}
\par\smallskip\noindent\begin{tabularx}{\linewidth}{@{}*{2}{L}@{}}
\optlabel{A}\,P$\to$2, Q$\to$3, R$\to$4, S$\to$1 & \optlabel{B}\,P$\to$2, Q$\to$3, R$\to$1, S$\to$4 \\[3pt]
\optlabel{C}\,P$\to$2, Q$\to$4, R$\to$3, S$\to$1 & \optlabel{D}\,P$\to$4, Q$\to$3, R$\to$2, S$\to$1
\end{tabularx}
\end{qitem}
\qsep

\needspace{7\baselineskip}
\begin{qitem}{16}
Match each classical algorithm (List-I) with its running time for an $n$-bit input $N$ (List-II).
\par\smallskip\noindent{\renewcommand{\arraystretch}{1.35}%
\begin{tabularx}{\linewidth}{@{}c>{\raggedright\arraybackslash}p{0.33\linewidth}@{\hspace{1.2em}}cL@{}}
\toprule
\multicolumn{2}{@{}l}{\textbf{List-I}} & \multicolumn{2}{l@{}}{\textbf{List-II}} \\ \midrule
\textbf{(P)} & trial division & \textbf{(1)} & $\exp\big((64/9)^{1/3}(\ln N)^{1/3}(\ln\ln N)^{2/3}(1 + o(1))\big)$ (sub-exponential) \\
\textbf{(Q)} & Pollard's rho & \textbf{(2)} & $O(N^{1/4})$ expected, i.e.\ $O(2^{n/4})$ (heuristic) \\
\textbf{(R)} & general number field sieve & \textbf{(3)} & $O(n)$ division steps (polynomial) \\
\textbf{(S)} & Euclidean $\gcd(a, N)$ & \textbf{(4)} & $O(\sqrt N) = O(2^{n/2})$ divisions \\
\bottomrule
\end{tabularx}}
\par\smallskip\noindent\begin{tabularx}{\linewidth}{@{}*{2}{L}@{}}
\optlabel{A}\,P$\to$4, Q$\to$2, R$\to$3, S$\to$1 & \optlabel{B}\,P$\to$4, Q$\to$2, R$\to$1, S$\to$3 \\[3pt]
\optlabel{C}\,P$\to$2, Q$\to$4, R$\to$1, S$\to$3 & \optlabel{D}\,P$\to$4, Q$\to$1, R$\to$2, S$\to$3
\end{tabularx}
\end{qitem}
\qsep

\needspace{7\baselineskip}
\begin{qitem}{17}
\textbf{Assertion (A):} If $x^2 \equiv 1 \pmod N$ but $x \not\equiv \pm1 \pmod N$, then $\gcd(x - 1, N)$ is a non-trivial factor of $N$.\\[1mm]
\textbf{Reason (R):} In that case $N$ divides $(x - 1)(x + 1)$ but divides neither $x - 1$ nor $x + 1$.
\begin{itemize}[leftmargin=2.6em, labelsep=0.4em, itemsep=2pt, topsep=3pt]
\item[\optlabel{A}] Both Assertion (A) and Reason (R) are true, and R is the correct explanation of A
\item[\optlabel{B}] Both A and R are true, but R is \emph{not} the correct explanation of A
\item[\optlabel{C}] A is true, but R is false
\item[\optlabel{D}] A is false, but R is true
\end{itemize}
\end{qitem}
\qsep

\needspace{7\baselineskip}
\begin{qitem}{18}
\textbf{Statement I:} If $\gcd(a, N) = 1$, then $a^{\varphi(N)} \equiv 1 \pmod N$.\\[1mm]
\textbf{Statement II:} The multiplicative order of every $a$ coprime to $N$ equals $\varphi(N)$.
\begin{itemize}[leftmargin=2.6em, labelsep=0.4em, itemsep=2pt, topsep=3pt]
\item[\optlabel{A}] Both Statement I and Statement II are true
\item[\optlabel{B}] Both Statement I and Statement II are false
\item[\optlabel{C}] Statement I is true, but Statement II is false
\item[\optlabel{D}] Statement I is false, but Statement II is true
\end{itemize}
\end{qitem}
\qsep

\needspace{20\baselineskip}
\begin{passagebox}{Paragraph for Questions 19 and 20}
Shor's reduction is applied by hand to $N = 91 = 7\times13$ with the base $a = 3$ (note $\gcd(3, 91) = 1$). Recall that the
order $r$ of $a$ modulo $N$ is the least $r > 0$ with $a^r \equiv 1 \pmod N$, and that if $r$ is even and $a^{r/2} \not\equiv -1 \pmod N$,
then $\gcd(a^{r/2} \pm 1, N)$ are non-trivial factors.
\end{passagebox}
\needspace{7\baselineskip}
\begin{qitem}{19}
The order $r$ of $3$ modulo $91$ is
\par\smallskip\noindent\begin{tabularx}{\linewidth}{@{}*{4}{L}@{}}
\optlabel{A}\,$6$ & \optlabel{B}\,$3$ & \optlabel{C}\,$72$ & \optlabel{D}\,$12$
\end{tabularx}
\end{qitem}
\qsep

\needspace{7\baselineskip}
\begin{qitem}{20}
The values of $\gcd(3^{r/2} - 1, 91)$ and $\gcd(3^{r/2} + 1, 91)$ are, respectively,
\par\smallskip\noindent\begin{tabularx}{\linewidth}{@{}*{4}{L}@{}}
\optlabel{A}\,$13$ and $7$ & \optlabel{B}\,$7$ and $13$ & \optlabel{C}\,$1$ and $91$ & \optlabel{D}\,$26$ and $28$
\end{tabularx}
\end{qitem}
\qsep

\newpage

\begin{sectionbanner}SECTION B \enspace$\cdot$\enspace Quantum Computing Foundations \hfill Q.21--40\end{sectionbanner}

{\small \textbf{Syllabus:} registers and superposition, entanglement, unitary modular multiplication, eigenvalues and eigenvectors, phase kickback, Hadamard, phase and controlled gates, interference, reversibility, measurement, simulation memory and phase estimation.\\ \textbf{Formats:} Single-correct MCQ (Q.21--34); Matching lists (Q.35--36); Assertion--Reason (Q.37); Statement I / II (Q.38); Paragraph based (Q.39--40). Each question: $+3$ / $0$ / $-1$.}
\par\medskip

\needspace{7\baselineskip}
\begin{qitem}{21}
After $H^{\otimes8}$ acts on $\ket{0}^{\otimes8}$ (the counting register of the notebook for $N = 15$), the amplitude of each basis state $\ket x$ is
\par\smallskip\noindent\begin{tabularx}{\linewidth}{@{}*{4}{L}@{}}
\optlabel{A}\,$1/\sqrt8$ & \optlabel{B}\,$1/16$ & \optlabel{C}\,$1/8$ & \optlabel{D}\,$1/256$
\end{tabularx}
\end{qitem}
\qsep

\needspace{7\baselineskip}
\begin{qitem}{22}
The state $\frac{1}{\sqrt8}\sum_{x=0}^{7}\ket x\ket{2^x \bmod 15}$ is prepared and the \emph{second} register is measured with result $4$. The probability
of this result and the post-measurement state of the first register are
\par\smallskip\noindent\begin{tabularx}{\linewidth}{@{}*{2}{L}@{}}
\optlabel{A}\,$1/4$; $\ \frac{1}{\sqrt2}(\ket2 - \ket6)$ & \optlabel{B}\,$1/4$; $\ \frac{1}{\sqrt2}(\ket2 + \ket6)$ \\[3pt]
\optlabel{C}\,$1/2$; $\ \frac{1}{\sqrt2}(\ket2 + \ket6)$ & \optlabel{D}\,$1/8$; $\ \ket2$
\end{tabularx}
\end{qitem}
\qsep

\needspace{7\baselineskip}
\begin{qitem}{23}
Why does the operator $U_a\ket y = \ket{ay \bmod N}$ (for $0 \le y < N$; identity for $y \ge N$) require $\gcd(a, N) = 1$ to be a valid quantum gate?
\begin{itemize}[leftmargin=2.6em, labelsep=0.4em, itemsep=2pt, topsep=3pt]
\item[\optlabel{A}] because otherwise $U_a$ would not be Hermitian
\item[\optlabel{B}] because the eigenvalues of $U_a$ must be real
\item[\optlabel{C}] because a quantum gate must map $\ket0$ to $\ket0$
\item[\optlabel{D}] only then is $y \mapsto ay \bmod N$ a permutation of $\{0, \dots, N-1\}$, so $U_a$ is a permutation matrix and hence unitary
\end{itemize}
\end{qitem}
\qsep

\needspace{7\baselineskip}
\begin{qitem}{24}
On the subspace spanned by $\{\ket1, \ket2, \ket4, \ket8\}$, the operator $U_2\ket y = \ket{2y \bmod 15}$ has the eigenvalues
\par\smallskip\noindent\begin{tabularx}{\linewidth}{@{}*{2}{L}@{}}
\optlabel{A}\,$1, 2, 4, 8$ & \optlabel{B}\,$e^{2\pi i/15}, e^{4\pi i/15}, e^{8\pi i/15}, e^{16\pi i/15}$ \\[3pt]
\optlabel{C}\,$1, -1$ (each twice) & \optlabel{D}\,$1, i, -1, -i$
\end{tabularx}
\end{qitem}
\qsep

\needspace{7\baselineskip}
\begin{qitem}{25}
With $\ket{u_s} = \frac{1}{\sqrt r}\sum_{k=0}^{r-1}e^{-2\pi isk/r}\ket{a^k \bmod N}$, the superposition $\frac{1}{\sqrt r}\sum_{s=0}^{r-1}\ket{u_s}$ equals
\par\smallskip\noindent\begin{tabularx}{\linewidth}{@{}*{2}{L}@{}}
\optlabel{A}\,$\ket0$ & \optlabel{B}\,$\ket a$ \\[3pt]
\optlabel{C}\,$\ket1$ & \optlabel{D}\,$\frac{1}{\sqrt N}\sum_{y=0}^{N-1}\ket y$
\end{tabularx}
\end{qitem}
\qsep

\needspace{7\baselineskip}
\begin{qitem}{26}
A control qubit in $\ket+$ controls a unitary $U$ acting on its eigenvector $\ket u$, $U\ket u = e^{i\phi}\ket u$, with $\phi = 2\pi/3$. Afterwards, the
expectation value $\langle X\rangle$ of the control qubit is
\par\smallskip\noindent\begin{tabularx}{\linewidth}{@{}*{4}{L}@{}}
\optlabel{A}\,$1/2$ & \optlabel{B}\,$1$ & \optlabel{C}\,$-1/2$ & \optlabel{D}\,$-\sqrt3/2$
\end{tabularx}
\end{qitem}
\qsep

\needspace{7\baselineskip}
\begin{qitem}{27}
In phase estimation, counting qubit $j$ controls $U^{2^j}$. If $U\ket u = e^{2\pi i/6}\ket u$, the phase kicked back onto counting qubit $j = 2$ is
\par\smallskip\noindent\begin{tabularx}{\linewidth}{@{}*{4}{L}@{}}
\optlabel{A}\,$2\pi/3$ & \optlabel{B}\,$4\pi/3$ & \optlabel{C}\,$\pi/3$ & \optlabel{D}\,$2\pi$
\end{tabularx}
\end{qitem}
\qsep

\needspace{7\baselineskip}
\begin{qitem}{28}
$(H\otimes H)\ket{11}$ equals
\par\smallskip\noindent\begin{tabularx}{\linewidth}{@{}*{2}{L}@{}}
\optlabel{A}\,$\frac12(\ket{00} + \ket{01} + \ket{10} + \ket{11})$ & \optlabel{B}\,$\frac12(\ket{00} - \ket{01} - \ket{10} - \ket{11})$ \\[3pt]
\optlabel{C}\,$\frac12(\ket{00} - \ket{01} - \ket{10} + \ket{11})$ & \optlabel{D}\,$\frac{1}{\sqrt2}(\ket{00} - \ket{11})$
\end{tabularx}
\end{qitem}
\qsep

\needspace{7\baselineskip}
\begin{qitem}{29}
Two computational paths lead to the same outcome with amplitudes $\frac12$ and $\frac12e^{i\theta}$. For $\theta = 2\pi/3$ the probability of that outcome is
\par\smallskip\noindent\begin{tabularx}{\linewidth}{@{}*{4}{L}@{}}
\optlabel{A}\,$0$ & \optlabel{B}\,$1/4$ & \optlabel{C}\,$3/4$ & \optlabel{D}\,$1/2$
\end{tabularx}
\end{qitem}
\qsep

\needspace{7\baselineskip}
\begin{qitem}{30}
Why does Shor's circuit compute $\ket x\ket1 \mapsto \ket x\ket{a^x \bmod N}$ instead of $\ket x \mapsto \ket{a^x \bmod N}$ in place?
\begin{itemize}[leftmargin=2.6em, labelsep=0.4em, itemsep=2pt, topsep=3pt]
\item[\optlabel{A}] the map $x \mapsto a^x \bmod N$ is not injective ($x$ and $x + r$ give the same value), so it cannot be a unitary on a single register; keeping $x$ makes it reversible
\item[\optlabel{B}] because measurement requires two registers
\item[\optlabel{C}] because the Hadamard gates must act on a different register than the QFT
\item[\optlabel{D}] because a single register cannot hold numbers larger than $N$
\end{itemize}
\end{qitem}
\qsep

\needspace{7\baselineskip}
\begin{qitem}{31}
A controlled-SWAP (Fredkin) gate is compiled as $\mathrm{CNOT}_{b\to a}\cdot\mathrm{Toffoli}_{c,a\to b}\cdot\mathrm{CNOT}_{b\to a}$, and each Toffoli uses the
standard decomposition with 6 CNOTs. The number of CNOTs per controlled-SWAP is
\par\smallskip\noindent\begin{tabularx}{\linewidth}{@{}*{4}{L}@{}}
\optlabel{A}\,$8$ & \optlabel{B}\,$6$ & \optlabel{C}\,$3$ & \optlabel{D}\,$9$
\end{tabularx}
\end{qitem}
\qsep

\needspace{7\baselineskip}
\begin{qitem}{32}
\texttt{Shor.ipynb} measures only the counting register. Compared with also measuring the work register first, the distribution of the counting
outcome $y$ is
\begin{itemize}[leftmargin=2.6em, labelsep=0.4em, itemsep=2pt, topsep=3pt]
\item[\optlabel{A}] different: without measuring the work register no periodic structure appears
\item[\optlabel{B}] different: the peaks become twice as wide
\item[\optlabel{C}] undefined, because entangled registers cannot be measured separately
\item[\optlabel{D}] exactly the same: a measurement of the other register that is never used cannot change the marginal statistics of $y$
\end{itemize}
\end{qitem}
\qsep

\needspace{7\baselineskip}
\begin{qitem}{33}
In \texttt{Shor.ipynb} the work register has $n = $ \texttt{N.bit\_length()} qubits and the counting register $t = 2n$. The total number of qubits for
$N = 143$ is
\par\smallskip\noindent\begin{tabularx}{\linewidth}{@{}*{4}{L}@{}}
\optlabel{A}\,$24$ & \optlabel{B}\,$21$ & \optlabel{C}\,$16$ & \optlabel{D}\,$22$
\end{tabularx}
\end{qitem}
\qsep

\needspace{7\baselineskip}
\begin{qitem}{34}
The memory of a double-precision state vector (16 bytes per amplitude) for the 24-qubit circuit of the previous question is about
\par\smallskip\noindent\begin{tabularx}{\linewidth}{@{}*{4}{L}@{}}
\optlabel{A}\,$268$ MB & \optlabel{B}\,$4.3$ GB & \optlabel{C}\,$384$ bytes & \optlabel{D}\,$16.8$ MB
\end{tabularx}
\end{qitem}
\qsep

\needspace{7\baselineskip}
\begin{qitem}{35}
Match each operation of the order-finding circuit (List-I) with its role (List-II).
\par\smallskip\noindent{\renewcommand{\arraystretch}{1.35}%
\begin{tabularx}{\linewidth}{@{}c>{\raggedright\arraybackslash}p{0.33\linewidth}@{\hspace{1.2em}}cL@{}}
\toprule
\multicolumn{2}{@{}l}{\textbf{List-I}} & \multicolumn{2}{l@{}}{\textbf{List-II}} \\ \midrule
\textbf{(P)} & $H$ on every counting qubit & \textbf{(1)} & turns the phases $2\pi\cdot2^js/r$ into a binary estimate $y \approx 2^ts/r$ \\
\textbf{(Q)} & controlled $U_{a^{2^j} \bmod N}$ & \textbf{(2)} & prepares the work register in $\ket1$, an equal superposition of the eigenvectors $\ket{u_s}$ \\
\textbf{(R)} & inverse QFT on the counting register & \textbf{(3)} & creates the uniform superposition over all exponents $x$ \\
\textbf{(S)} & $X$ on the last work qubit & \textbf{(4)} & performs modular exponentiation / kicks the eigenphases back onto the counting qubits \\
\bottomrule
\end{tabularx}}
\par\smallskip\noindent\begin{tabularx}{\linewidth}{@{}*{2}{L}@{}}
\optlabel{A}\,P$\to$3, Q$\to$4, R$\to$2, S$\to$1 & \optlabel{B}\,P$\to$4, Q$\to$3, R$\to$1, S$\to$2 \\[3pt]
\optlabel{C}\,P$\to$3, Q$\to$4, R$\to$1, S$\to$2 & \optlabel{D}\,P$\to$3, Q$\to$1, R$\to$4, S$\to$2
\end{tabularx}
\end{qitem}
\qsep

\needspace{7\baselineskip}
\begin{qitem}{36}
The phase gate is $R_1(\lambda) = \operatorname{diag}(1, e^{i\lambda})$ (controlled versions appear in the QFT). Match each angle (List-I) with the named gate (List-II).
\par\smallskip\noindent{\renewcommand{\arraystretch}{1.35}%
\begin{tabularx}{\linewidth}{@{}c>{\raggedright\arraybackslash}p{0.33\linewidth}@{\hspace{1.2em}}cL@{}}
\toprule
\multicolumn{2}{@{}l}{\textbf{List-I}} & \multicolumn{2}{l@{}}{\textbf{List-II}} \\ \midrule
\textbf{(P)} & $R_1(\pi)$ & \textbf{(1)} & $S^\dagger$ \\
\textbf{(Q)} & $R_1(\pi/2)$ & \textbf{(2)} & $T$ \\
\textbf{(R)} & $R_1(\pi/4)$ & \textbf{(3)} & $Z$ \\
\textbf{(S)} & $R_1(-\pi/2)$ & \textbf{(4)} & $S$ \\
\bottomrule
\end{tabularx}}
\par\smallskip\noindent\begin{tabularx}{\linewidth}{@{}*{2}{L}@{}}
\optlabel{A}\,P$\to$3, Q$\to$2, R$\to$4, S$\to$1 & \optlabel{B}\,P$\to$4, Q$\to$3, R$\to$2, S$\to$1 \\[3pt]
\optlabel{C}\,P$\to$3, Q$\to$4, R$\to$1, S$\to$2 & \optlabel{D}\,P$\to$3, Q$\to$4, R$\to$2, S$\to$1
\end{tabularx}
\end{qitem}
\qsep

\needspace{7\baselineskip}
\begin{qitem}{37}
\textbf{Assertion (A):} When the target register is in an eigenvector of $U$, a controlled-$U$ gate leaves the target register unchanged.\\[1mm]
\textbf{Reason (R):} Controlled gates never change the state of their control qubit.
\begin{itemize}[leftmargin=2.6em, labelsep=0.4em, itemsep=2pt, topsep=3pt]
\item[\optlabel{A}] Both Assertion (A) and Reason (R) are true, and R is the correct explanation of A
\item[\optlabel{B}] Both A and R are true, but R is \emph{not} the correct explanation of A
\item[\optlabel{C}] A is true, but R is false
\item[\optlabel{D}] A is false, but R is true
\end{itemize}
\end{qitem}
\qsep

\needspace{7\baselineskip}
\begin{qitem}{38}
\textbf{Statement I:} Because the circuit evaluates $a^x \bmod N$ for all $x$ in superposition, measuring the registers directly reveals the period $r$.\\[1mm]
\textbf{Statement II:} Interference produced by the inverse QFT is needed to turn the periodic structure into high-probability measurement outcomes.
\begin{itemize}[leftmargin=2.6em, labelsep=0.4em, itemsep=2pt, topsep=3pt]
\item[\optlabel{A}] Both Statement I and Statement II are true
\item[\optlabel{B}] Both Statement I and Statement II are false
\item[\optlabel{C}] Statement I is true, but Statement II is false
\item[\optlabel{D}] Statement I is false, but Statement II is true
\end{itemize}
\end{qitem}
\qsep

\needspace{20\baselineskip}
\begin{passagebox}{Paragraph for Questions 39 and 40}
Quantum phase estimation (QPE) with $t = 3$ counting qubits is applied to a unitary $U$ and an eigenvector $\ket u$ with
$U\ket u = e^{2\pi i\varphi}\ket u$, $0 \le \varphi < 1$. After the Hadamards, the controlled-$U^{2^j}$ gates and the inverse QFT, the counting register is
measured and gives $y \in \{0, \dots, 7\}$ with probability
\[ P(y) = \Big|\frac18\sum_{x=0}^{7}e^{2\pi ix(\varphi - y/8)}\Big|^2 = \frac{\sin^2\big(8\pi(\varphi - y/8)\big)}{64\,\sin^2\big(\pi(\varphi - y/8)\big)}\quad(P = 1\text{ if }\varphi = y/8). \]
The estimate of the phase is $y/8$; $y$ is written in binary with the most significant bit first.
\end{passagebox}
\needspace{7\baselineskip}
\begin{qitem}{39}
For $\varphi = 3/8$ the measurement gives
\par\smallskip\noindent\begin{tabularx}{\linewidth}{@{}*{2}{L}@{}}
\optlabel{A}\,$y = 3$ (bits $011$) with probability $1$ & \optlabel{B}\,$y = 6$ (bits $110$) with probability $1$ \\[3pt]
\optlabel{C}\,$y = 3$ or $y = 4$ with probability $1/2$ each & \optlabel{D}\,$y = 3$ with probability $1/2$
\end{tabularx}
\end{qitem}
\qsep

\needspace{7\baselineskip}
\begin{qitem}{40}
For $\varphi = 1/3$ the most likely outcome and its probability are
\par\smallskip\noindent\begin{tabularx}{\linewidth}{@{}*{2}{L}@{}}
\optlabel{A}\,$y = 3$ with $P \approx 0.41$ & \optlabel{B}\,$y = 3$ with $P \approx 0.69$ \\[3pt]
\optlabel{C}\,$y = 3$ with $P = 1$ & \optlabel{D}\,$y = 2$ or $3$ with $P = 1/2$ each
\end{tabularx}
\end{qitem}
\qsep

\newpage

\begin{sectionbanner}SECTION C \enspace$\cdot$\enspace Shor's Algorithm Fundamentals \hfill Q.41--60\end{sectionbanner}

{\small \textbf{Syllabus:} the steps of the algorithm, choice of the base, register sizes, outcome distributions, continued fractions, success and failure cases, repetition, classical post-processing and the functions of \texttt{Shor.ipynb}.\\ \textbf{Formats:} Single-correct MCQ (Q.41--54); Matching lists (Q.55--56); Assertion--Reason (Q.57); Statement I / II (Q.58); Paragraph based (Q.59--60). Each question: $+3$ / $0$ / $-1$.}
\par\medskip

\needspace{7\baselineskip}
\begin{qitem}{41}
Which sequence describes Shor's factoring algorithm for an odd composite $N$ that is not a prime power?
\begin{itemize}[leftmargin=2.6em, labelsep=0.4em, itemsep=2pt, topsep=3pt]
\item[\optlabel{A}] pick random $a$; if $\gcd(a, N) > 1$ stop (factor found); find the order $r$ of $a$ by quantum phase estimation and continued fractions; if $r$ is even and $a^{r/2} \not\equiv -1$, output $\gcd(a^{r/2} \pm 1, N)$; otherwise repeat
\item[\optlabel{B}] pick random $a$; measure $a^x \bmod N$ in superposition; the most frequent value is a factor
\item[\optlabel{C}] find the order $r$ by quantum phase estimation; the factors are $r$ and $N/r$
\item[\optlabel{D}] apply the QFT to $\ket N$; measure; the outcome is a factor of $N$
\end{itemize}
\end{qitem}
\qsep

\needspace{7\baselineskip}
\begin{qitem}{42}
For $N = 15$ and $a = 2$, the order $r$ and the factors obtained from $\gcd(a^{r/2} \pm 1, N)$ are
\par\smallskip\noindent\begin{tabularx}{\linewidth}{@{}*{2}{L}@{}}
\optlabel{A}\,$r = 8$; $\ 3$ and $5$ & \optlabel{B}\,$r = 2$; $\ 3$ and $5$ \\[3pt]
\optlabel{C}\,$r = 4$; $\ 1$ and $15$ & \optlabel{D}\,$r = 4$; $\ \gcd(3, 15) = 3$, $\gcd(5, 15) = 5$
\end{tabularx}
\end{qitem}
\qsep

\needspace{7\baselineskip}
\begin{qitem}{43}
For $N = 15$ and $a = 14$, the order finding succeeds perfectly. The classical post-processing then
\begin{itemize}[leftmargin=2.6em, labelsep=0.4em, itemsep=2pt, topsep=3pt]
\item[\optlabel{A}] succeeds with factors $3$ and $5$
\item[\optlabel{B}] fails because $r$ is odd
\item[\optlabel{C}] fails: $r = 2$ but $14^{1} \equiv -1 \pmod{15}$, so $\gcd(14 \pm 1, 15)$ gives only the trivial factors $15$ and $1$; another $a$ is needed
\item[\optlabel{D}] fails because $\gcd(14, 15) = 1$
\end{itemize}
\end{qitem}
\qsep

\needspace{7\baselineskip}
\begin{qitem}{44}
For $N = 21$ and $a = 4$, the outcome of the classical checks is
\begin{itemize}[leftmargin=2.6em, labelsep=0.4em, itemsep=2pt, topsep=3pt]
\item[\optlabel{A}] $r = 6$, factors $3$ and $7$
\item[\optlabel{B}] $r = 3$, factors $\gcd(4^1 - 1, 21) = 3$ and $\gcd(4^1 + 1, 21) = 1$
\item[\optlabel{C}] $r = 3$ is odd, so $a^{r/2}$ is not an integer power; the attempt fails and a new $a$ is chosen
\item[\optlabel{D}] $\gcd(4, 21) = 7$, a lucky factor
\end{itemize}
\end{qitem}
\qsep

\needspace{7\baselineskip}
\begin{qitem}{45}
For $N = 21$ the randomly chosen base is $a = 6$. What happens?
\begin{itemize}[leftmargin=2.6em, labelsep=0.4em, itemsep=2pt, topsep=3pt]
\item[\optlabel{A}] order finding gives $r = 6$ and the factors $3$, $7$
\item[\optlabel{B}] $U_6$ is a valid unitary and phase estimation returns $r = 2$
\item[\optlabel{C}] the algorithm fails because $a$ is even
\item[\optlabel{D}] $\gcd(6, 21) = 3$ is already a non-trivial factor; no quantum computation is needed
\end{itemize}
\end{qitem}
\qsep

\needspace{7\baselineskip}
\begin{qitem}{46}
With the register sizes of \texttt{Shor.ipynb} ($n = $ \texttt{N.bit\_length()}, $t = 2n$), the circuit for $N = 33$ uses
\par\smallskip\noindent\begin{tabularx}{\linewidth}{@{}*{2}{L}@{}}
\optlabel{A}\,$33$ qubits & \optlabel{B}\,$15$ qubits ($n = 5$, $t = 10$) \\[3pt]
\optlabel{C}\,$12$ qubits & \optlabel{D}\,$18$ qubits ($n = 6$, $t = 12$)
\end{tabularx}
\end{qitem}
\qsep

\needspace{7\baselineskip}
\begin{qitem}{47}
For $N = 15$, $a = 11$ and $t = 8$ counting qubits, the ideal outcome distribution and the result of post-processing $y = 128$ are
\begin{itemize}[leftmargin=2.6em, labelsep=0.4em, itemsep=2pt, topsep=3pt]
\item[\optlabel{A}] a uniform distribution, because $11$ is prime
\item[\optlabel{B}] peaks at $y \in \{0, 128\}$ with probability $1/2$ each; $128/256 = 1/2 \Rightarrow r = 2$, $\gcd(10, 15) = 5$, $\gcd(12, 15) = 3$
\item[\optlabel{C}] peaks at $y \in \{0, 128\}$; $r = 2$, but $11 \equiv -1$, so the factors are trivial
\item[\optlabel{D}] peaks at $y \in \{0, 64, 128, 192\}$; $r = 4$
\end{itemize}
\end{qitem}
\qsep

\needspace{7\baselineskip}
\begin{qitem}{48}
For $N = 33$, $a = 5$ and $t = 12$, the measurement gives $y = 1229$. The continued-fraction step and the factors are
\begin{itemize}[leftmargin=2.6em, labelsep=0.4em, itemsep=2pt, topsep=3pt]
\item[\optlabel{A}] $1229/4096 \approx 3/10 \Rightarrow r = 10$; $5^5 \equiv -1$, trivial factors
\item[\optlabel{B}] $1229/4096 \approx 3/10 \Rightarrow r = 10$; $5^5 \equiv 23$, $\gcd(22, 33) = 11$, $\gcd(24, 33) = 3$
\item[\optlabel{C}] $1229/4096 \approx 1/3 \Rightarrow r = 3$; the attempt fails
\item[\optlabel{D}] $1229/4096 \approx 3/10 \Rightarrow r = 3$; factors $3$ and $11$
\end{itemize}
\end{qitem}
\qsep

\needspace{7\baselineskip}
\begin{qitem}{49}
For $N = 33$, $a = 5$, $t = 12$, the outcome is $y = 2048$. The procedure of \texttt{candidate\_orders} returns
\begin{itemize}[leftmargin=2.6em, labelsep=0.4em, itemsep=2pt, topsep=3pt]
\item[\optlabel{A}] the order $2$, because the denominator is always the order
\item[\optlabel{B}] the order $5$
\item[\optlabel{C}] no candidate, because $y = 2^{t-1}$ is never useful
\item[\optlabel{D}] the fraction $1/2$ and, after testing multiples of $2$, the order $10$ (since $5^2, 5^4, 5^6, 5^8 \not\equiv 1$, $5^{10} \equiv 1$)
\end{itemize}
\end{qitem}
\qsep

\needspace{7\baselineskip}
\begin{qitem}{50}
The textbook condition $2^t \ge N^2$ guarantees that continued fractions recover $s/r$ from the nearest outcome. For $N = 33$, the smallest such
$t$ and the value used by \texttt{Shor.ipynb} are
\par\smallskip\noindent\begin{tabularx}{\linewidth}{@{}*{2}{L}@{}}
\optlabel{A}\,$t = 11$; the notebook uses $t = 2n = 12$ & \optlabel{B}\,$t = 10$; the notebook uses $12$ \\[3pt]
\optlabel{C}\,$t = 6$; the notebook uses $12$ & \optlabel{D}\,$t = 12$; the notebook uses $11$
\end{tabularx}
\end{qitem}
\qsep

\needspace{7\baselineskip}
\begin{qitem}{51}
For the notebook's default ($N = 15$, $a = 7$, $t = 8$) the ideal outcomes are $y \in \{0, 64, 128, 192\}$, each with probability $1/4$. Using the
multiples test, the exact probability that one shot yields the factors is
\par\smallskip\noindent\begin{tabularx}{\linewidth}{@{}*{2}{L}@{}}
\optlabel{A}\,$1$ & \optlabel{B}\,$3/4$ --- only $y = 0$ fails \\[3pt]
\optlabel{C}\,$1/2$ --- only $y = 64$ and $192$ succeed & \optlabel{D}\,$1/4$ --- only $y = 64$ succeeds
\end{tabularx}
\end{qitem}
\qsep

\needspace{7\baselineskip}
\begin{qitem}{52}
Suppose $r = 12$ divides $2^t$, so every outcome is exactly $y = s2^t/r$ with $s$ uniform in $\{0, \dots, 11\}$. \emph{Without} the multiples test,
the fraction of shots whose reduced denominator equals $r$ is
\par\smallskip\noindent\begin{tabularx}{\linewidth}{@{}*{4}{L}@{}}
\optlabel{A}\,$\varphi(12)/12 = 1/3$ & \optlabel{B}\,$1/12$ & \optlabel{C}\,$1/2$ & \optlabel{D}\,$11/12$
\end{tabularx}
\end{qitem}
\qsep

\needspace{7\baselineskip}
\begin{qitem}{53}
What is the role of the (inverse) quantum Fourier transform \emph{within} Shor's algorithm?
\begin{itemize}[leftmargin=2.6em, labelsep=0.4em, itemsep=2pt, topsep=3pt]
\item[\optlabel{A}] it is a subroutine of phase estimation that turns the periodic phase pattern on the counting register into peaks near $y = k2^t/r$; it does not factor $N$ by itself
\item[\optlabel{B}] it is the whole of Shor's algorithm; the remaining steps are optional
\item[\optlabel{C}] it computes $a^x \bmod N$
\item[\optlabel{D}] it finds the factors $p$ and $q$ directly from $\ket N$
\end{itemize}
\end{qitem}
\qsep

\needspace{7\baselineskip}
\begin{qitem}{54}
\texttt{build\_shor\_kernel} computes $c_j = a^{2^j} \bmod N$ classically and skips $c_j = 1$. For $N = 15$, $a = 7$, $t = 8$, the number of controlled
multiplications actually applied is
\par\smallskip\noindent\begin{tabularx}{\linewidth}{@{}*{2}{L}@{}}
\optlabel{A}\,$4$ & \optlabel{B}\,$2$ ($c_0 = 7$, $c_1 = 4$; $c_j = 1$ for $j \ge 2$) \\[3pt]
\optlabel{C}\,$1$ & \optlabel{D}\,$8$
\end{tabularx}
\end{qitem}
\qsep

\needspace{7\baselineskip}
\begin{qitem}{55}
Match each function of \texttt{Shor.ipynb} (List-I) with what it does (List-II).
\par\smallskip\noindent{\renewcommand{\arraystretch}{1.35}%
\begin{tabularx}{\linewidth}{@{}c>{\raggedright\arraybackslash}p{0.33\linewidth}@{\hspace{1.2em}}cL@{}}
\toprule
\multicolumn{2}{@{}l}{\textbf{List-I}} & \multicolumn{2}{l@{}}{\textbf{List-II}} \\ \midrule
\textbf{(P)} & \texttt{mulmod\_matrix} & \textbf{(1)} & brute-force order of $a$ modulo $m$ as a classical reference (feasible only for tiny $m$) \\
\textbf{(Q)} & \texttt{build\_shor\_kernel} & \textbf{(2)} & continued-fraction denominator of $y/2^t$ and the smallest multiple $k r'$ with $a^{kr'} \equiv 1$ \\
\textbf{(R)} & \texttt{candidate\_orders} & \textbf{(3)} & builds the CUDA-Q circuit: Hadamards, controlled $U_{a^{2^j}}$, inverse QFT, measurement \\
\textbf{(S)} & \texttt{classical\_order} & \textbf{(4)} & permutation matrix $\ket y \mapsto \ket{cy \bmod N}$ for $y < N$, identity above \\
\bottomrule
\end{tabularx}}
\par\smallskip\noindent\begin{tabularx}{\linewidth}{@{}*{2}{L}@{}}
\optlabel{A}\,P$\to$4, Q$\to$2, R$\to$3, S$\to$1 & \optlabel{B}\,P$\to$4, Q$\to$3, R$\to$1, S$\to$2 \\[3pt]
\optlabel{C}\,P$\to$4, Q$\to$3, R$\to$2, S$\to$1 & \optlabel{D}\,P$\to$3, Q$\to$4, R$\to$2, S$\to$1
\end{tabularx}
\end{qitem}
\qsep

\needspace{7\baselineskip}
\begin{qitem}{56}
For $N = 33$, $a = 5$ ($r = 10$) and $t = 12$, match each outcome $y$ (List-I) with \texttt{Fraction(y, 4096).limit\_denominator(33)} (List-II).
\par\smallskip\noindent{\renewcommand{\arraystretch}{1.35}%
\begin{tabularx}{\linewidth}{@{}c>{\raggedright\arraybackslash}p{0.33\linewidth}@{\hspace{1.2em}}cL@{}}
\toprule
\multicolumn{2}{@{}l}{\textbf{List-I}} & \multicolumn{2}{l@{}}{\textbf{List-II}} \\ \midrule
\textbf{(P)} & $410$ & \textbf{(1)} & $2/5$ \\
\textbf{(Q)} & $819$ & \textbf{(2)} & $7/10$ \\
\textbf{(R)} & $1638$ & \textbf{(3)} & $1/10$ \\
\textbf{(S)} & $2867$ & \textbf{(4)} & $1/5$ \\
\bottomrule
\end{tabularx}}
\par\smallskip\noindent\begin{tabularx}{\linewidth}{@{}*{2}{L}@{}}
\optlabel{A}\,P$\to$4, Q$\to$3, R$\to$1, S$\to$2 & \optlabel{B}\,P$\to$3, Q$\to$4, R$\to$2, S$\to$1 \\[3pt]
\optlabel{C}\,P$\to$3, Q$\to$4, R$\to$1, S$\to$2 & \optlabel{D}\,P$\to$3, Q$\to$1, R$\to$4, S$\to$2
\end{tabularx}
\end{qitem}
\qsep

\needspace{7\baselineskip}
\begin{qitem}{57}
\textbf{Assertion (A):} Shor's order-finding reduction is the method of choice for factoring $N = 81 = 3^4$.\\[1mm]
\textbf{Reason (R):} For an odd prime power $N = p^k$, the only square roots of $1$ modulo $N$ are $\pm1$.
\begin{itemize}[leftmargin=2.6em, labelsep=0.4em, itemsep=2pt, topsep=3pt]
\item[\optlabel{A}] Both Assertion (A) and Reason (R) are true, and R is the correct explanation of A
\item[\optlabel{B}] Both A and R are true, but R is \emph{not} the correct explanation of A
\item[\optlabel{C}] A is true, but R is false
\item[\optlabel{D}] A is false, but R is true
\end{itemize}
\end{qitem}
\qsep

\needspace{7\baselineskip}
\begin{qitem}{58}
\textbf{Statement I:} A measured value $y = 0$ means that the order is $r = 1$.\\[1mm]
\textbf{Statement II:} When $r$ does not divide $2^t$, every measured $y$ yields the exact order through continued fractions.
\begin{itemize}[leftmargin=2.6em, labelsep=0.4em, itemsep=2pt, topsep=3pt]
\item[\optlabel{A}] Both Statement I and Statement II are true
\item[\optlabel{B}] Both Statement I and Statement II are false
\item[\optlabel{C}] Statement I is true, but Statement II is false
\item[\optlabel{D}] Statement I is false, but Statement II is true
\end{itemize}
\end{qitem}
\qsep

\needspace{20\baselineskip}
\begin{passagebox}{Paragraph for Questions 59 and 60}
Shor's algorithm is applied to $N = 39$ with base $a = 2$, using the register sizes of \texttt{Shor.ipynb}: the work register has
$n = $ \texttt{N.bit\_length()} $= 6$ qubits and the counting register $t = 2n = 12$ qubits, so the peaks of the outcome distribution lie near
$y \approx k\cdot4096/r$. A measured value $y$ is post-processed with \texttt{Fraction(y, 2**t).limit\_denominator(N)}; the denominator $r'$ of the
result, or the smallest multiple $kr'$ with $a^{kr'} \equiv 1 \pmod N$, is taken as the order. The powers of $2$ modulo $39$ are
$2, 4, 8, 16, 32, 25, 11, 22, 5, 10, 20, 1$.
\end{passagebox}
\needspace{7\baselineskip}
\begin{qitem}{59}
The order $r$ and the factors $\gcd(2^{r/2} - 1, 39)$, $\gcd(2^{r/2} + 1, 39)$ are
\par\smallskip\noindent\begin{tabularx}{\linewidth}{@{}*{2}{L}@{}}
\optlabel{A}\,$r = 12$; $\ 1$ and $39$ & \optlabel{B}\,$r = 12$; $\ 7$ and $13$ \\[3pt]
\optlabel{C}\,$r = 12$; $\ 3$ and $13$ & \optlabel{D}\,$r = 6$; $\ 3$ and $13$
\end{tabularx}
\end{qitem}
\qsep

\needspace{7\baselineskip}
\begin{qitem}{60}
The outcome $y = 1707$ is measured. The continued-fraction step gives
\par\smallskip\noindent\begin{tabularx}{\linewidth}{@{}*{2}{L}@{}}
\optlabel{A}\,$5/12$, so $r = 12$ directly & \optlabel{B}\,$1/2$, so the multiples test is needed \\[3pt]
\optlabel{C}\,$5/12$, so $r = 5$ & \optlabel{D}\,$2/5$, so $r = 5$
\end{tabularx}
\end{qitem}
\qsep

\newpage

\begin{sectionbanner}SECTION D \enspace$\cdot$\enspace Advanced Shor's Algorithm and Implementation \hfill Q.61--80\end{sectionbanner}

{\small \textbf{Syllabus:} resource estimation and published estimates, qubit-saving constructions, gate-level compilation, approximate QFT, noise, simulation limits, RSA and elliptic curves, post-quantum cryptography and fault tolerance.\\ \textbf{Formats:} Single-correct MCQ (Q.61--74); Matching lists (Q.75--76); Assertion--Reason (Q.77); Statement I / II (Q.78); Paragraph based (Q.79--80). Each question: $+3$ / $0$ / $-1$.}
\par\medskip

\needspace{7\baselineskip}
\begin{qitem}{61}
Beauregard's circuit for Shor's algorithm uses $2n + 3$ logical qubits. For RSA-2048 ($n = 2048$) this is
\par\smallskip\noindent\begin{tabularx}{\linewidth}{@{}*{4}{L}@{}}
\optlabel{A}\,$2051$ & \optlabel{B}\,$4096$ & \optlabel{C}\,$6147$ & \optlabel{D}\,$4099$
\end{tabularx}
\end{qitem}
\qsep

\needspace{7\baselineskip}
\begin{qitem}{62}
Which pair of published resource estimates for factoring RSA-2048 is correct?
\begin{itemize}[leftmargin=2.6em, labelsep=0.4em, itemsep=2pt, topsep=3pt]
\item[\optlabel{A}] both estimates state that RSA-2048 was factored on existing hardware
\item[\optlabel{B}] Gidney (2025): about $20$ billion qubits in about $8$ hours
\item[\optlabel{C}] Gidney and Ekerå (2019): about $20$ million noisy qubits in about $8$ hours; Gidney (2025): fewer than one million noisy qubits in less than a week
\item[\optlabel{D}] Gidney and Ekerå (2019): $4099$ physical qubits in one second
\end{itemize}
\end{qitem}
\qsep

\needspace{7\baselineskip}
\begin{qitem}{63}
With schoolbook arithmetic, modular exponentiation costs $O(n^3)$ gates. Doubling the bit length $n$ of $N$ multiplies the gate count by about
\par\smallskip\noindent\begin{tabularx}{\linewidth}{@{}*{4}{L}@{}}
\optlabel{A}\,$4$ & \optlabel{B}\,$2^{n}$ & \optlabel{C}\,$2$ & \optlabel{D}\,$8$
\end{tabularx}
\end{qitem}
\qsep

\needspace{7\baselineskip}
\begin{qitem}{64}
In the ``semiclassical'' (one-control-qubit) implementation of the inverse QFT used by Beauregard, the $t = 2n$ counting qubits are replaced by
\begin{itemize}[leftmargin=2.6em, labelsep=0.4em, itemsep=2pt, topsep=3pt]
\item[\optlabel{A}] no qubits: the QFT is computed classically before the circuit
\item[\optlabel{B}] $2n$ ancilla qubits that are never measured
\item[\optlabel{C}] a single qubit that is measured and reset $2n$ times; earlier measurement results classically control the phase corrections
\item[\optlabel{D}] $n$ qubits, because the QFT needs only half the bits
\end{itemize}
\end{qitem}
\qsep

\needspace{7\baselineskip}
\begin{qitem}{65}
For $N = 21$, $a = 2$, $t = 10$, an approximate inverse QFT keeps only controlled phases between qubits at distance $k - j \le m$. With $m = 2$ the
number of controlled phases and the (re-executed, exact) success probability are
\begin{itemize}[leftmargin=2.6em, labelsep=0.4em, itemsep=2pt, topsep=3pt]
\item[\optlabel{A}] $45$ phases; $0.817$
\item[\optlabel{B}] $2$ phases; $0.689$
\item[\optlabel{C}] $17$ phases; $0.817$ --- truncation never changes the result
\item[\optlabel{D}] $17$ phases; $0.689$ (full QFT: $45$ phases, $0.817$)
\end{itemize}
\end{qitem}
\qsep

\needspace{7\baselineskip}
\begin{qitem}{66}
The full inverse QFT on $t = 8$ counting qubits (notebook default for $N = 15$) contains how many controlled phase gates?
\par\smallskip\noindent\begin{tabularx}{\linewidth}{@{}*{4}{L}@{}}
\optlabel{A}\,$8$ & \optlabel{B}\,$64$ & \optlabel{C}\,$28$ & \optlabel{D}\,$56$
\end{tabularx}
\end{qitem}
\qsep

\needspace{7\baselineskip}
\begin{qitem}{67}
The recorded gate counts of \texttt{Shor.ipynb} for $N = 15$ are \texttt{x: 1, ccustom\_unitary\_constant: 2, swap: 4, h: 16, cr1: 28} (and
\texttt{mz: 8}). Why are there $16$ Hadamard gates?
\begin{itemize}[leftmargin=2.6em, labelsep=0.4em, itemsep=2pt, topsep=3pt]
\item[\optlabel{A}] each controlled multiplication contains $8$ Hadamards
\item[\optlabel{B}] Hadamards are inserted before every measurement of the work register
\item[\optlabel{C}] $16 = 2^4$ is the dimension of the work register
\item[\optlabel{D}] $8$ create the superposition on the counting qubits and $8$ belong to the inverse QFT
\end{itemize}
\end{qitem}
\qsep

\needspace{7\baselineskip}
\begin{qitem}{68}
Why does \texttt{Shor.ipynb} refuse to run on hardware targets?
\begin{itemize}[leftmargin=2.6em, labelsep=0.4em, itemsep=2pt, topsep=3pt]
\item[\optlabel{A}] the inverse QFT cannot be executed on hardware
\item[\optlabel{B}] Shor's algorithm is mathematically impossible on hardware
\item[\optlabel{C}] CUDA-Q has no hardware targets
\item[\optlabel{D}] its modular multipliers are exact custom unitaries (permutation matrices) registered with \texttt{cudaq.register\_operation}; hardware would need them decomposed into native gates, which the notebook does not do
\end{itemize}
\end{qitem}
\qsep

\needspace{7\baselineskip}
\begin{qitem}{69}
In the gate-level circuit, ``multiply by $7$ modulo $15$'' is built as three left rotations of the 4 work bits (each $3$ SWAPs) followed by $X$ on all
four bits. Its controlled version, with controlled SWAPs of $8$ CNOTs and controlled-$X$ $=$ CNOT, uses how many CNOTs?
\par\smallskip\noindent\begin{tabularx}{\linewidth}{@{}*{4}{L}@{}}
\optlabel{A}\,$9$ & \optlabel{B}\,$40$ & \optlabel{C}\,$76$ & \optlabel{D}\,$72$
\end{tabularx}
\end{qitem}
\qsep

\needspace{7\baselineskip}
\begin{qitem}{70}
In the gate-level circuit ($t = 3$, $a = 7$, $N = 15$) the counting outcome becomes uniformly random in the fully depolarized limit. The success
probability with the notebook's post-processing is then
\begin{itemize}[leftmargin=2.6em, labelsep=0.4em, itemsep=2pt, topsep=3pt]
\item[\optlabel{A}] $0$
\item[\optlabel{B}] $3/8$: $y = 2, 4, 6$ give $r = 4$ and the factors; $y = 0$ gives nothing; odd $y$ give $r' = 8$, and $7^4 \equiv 1$ gives only trivial factors
\item[\optlabel{C}] $3/4$, the same as without noise
\item[\optlabel{D}] $1/8$
\end{itemize}
\end{qitem}
\qsep

\needspace{7\baselineskip}
\begin{qitem}{71}
Using the notebook's $3n$-qubit construction, the state-vector memory (16 bytes per amplitude) for $N = 1023$ is about
\par\smallskip\noindent\begin{tabularx}{\linewidth}{@{}*{2}{L}@{}}
\optlabel{A}\,$268$ MB ($24$ qubits) & \optlabel{B}\,$17.2$ GB ($30$ qubits) \\[3pt]
\optlabel{C}\,$1.1$ TB ($36$ qubits) & \optlabel{D}\,$16.4$ kB
\end{tabularx}
\end{qitem}
\qsep

\needspace{7\baselineskip}
\begin{qitem}{72}
Toy RSA: $p = 5$, $q = 11$, $e = 7$. The private exponent $d$ and the ciphertext of $m = 2$ are
\par\smallskip\noindent\begin{tabularx}{\linewidth}{@{}*{4}{L}@{}}
\optlabel{A}\,$d = 23$; $c = 18$ & \optlabel{B}\,$d = 23$; $c = 14$ & \optlabel{C}\,$d = 3$; $c = 18$ & \optlabel{D}\,$d = 7$; $c = 2$
\end{tabularx}
\end{qitem}
\qsep

\needspace{7\baselineskip}
\begin{qitem}{73}
Which standard specifies NIST's lattice-based key-encapsulation mechanism intended to replace RSA/ECDH key exchange?
\par\smallskip\noindent\begin{tabularx}{\linewidth}{@{}*{2}{L}@{}}
\optlabel{A}\,FIPS 197 (AES) & \optlabel{B}\,FIPS 203 (ML-KEM) \\[3pt]
\optlabel{C}\,SP 800-90A (random bit generators) & \optlabel{D}\,FIPS 186 (DSA/ECDSA)
\end{tabularx}
\end{qitem}
\qsep

\needspace{7\baselineskip}
\begin{qitem}{74}
Which statement about the cryptographic impact of a large fault-tolerant quantum computer is correct?
\begin{itemize}[leftmargin=2.6em, labelsep=0.4em, itemsep=2pt, topsep=3pt]
\item[\optlabel{A}] Shor's algorithm breaks RSA and elliptic-curve cryptography; AES is affected only by Grover's quadratic speed-up, so AES-256 remains adequate
\item[\optlabel{B}] hash functions such as SHA-256 are broken by Shor's algorithm
\item[\optlabel{C}] Shor's algorithm breaks AES-256 in polynomial time
\item[\optlabel{D}] only RSA is affected; elliptic-curve cryptography is immune
\end{itemize}
\end{qitem}
\qsep

\needspace{7\baselineskip}
\begin{qitem}{75}
Match each fault-tolerance term (List-I) with its meaning (List-II).
\par\smallskip\noindent{\renewcommand{\arraystretch}{1.35}%
\begin{tabularx}{\linewidth}{@{}c>{\raggedright\arraybackslash}p{0.33\linewidth}@{\hspace{1.2em}}cL@{}}
\toprule
\multicolumn{2}{@{}l}{\textbf{List-I}} & \multicolumn{2}{l@{}}{\textbf{List-II}} \\ \midrule
\textbf{(P)} & logical qubit & \textbf{(1)} & minimum number of physical errors that can cause an undetected logical error \\
\textbf{(Q)} & physical qubit & \textbf{(2)} & a special resource state that enables non-Clifford gates such as $T$ or Toffoli in the surface code \\
\textbf{(R)} & code distance $d$ & \textbf{(3)} & an error-corrected qubit encoded in many noisy qubits \\
\textbf{(S)} & magic state & \textbf{(4)} & a noisy hardware qubit with error rates around $10^{-3}$ \\
\bottomrule
\end{tabularx}}
\par\smallskip\noindent\begin{tabularx}{\linewidth}{@{}*{2}{L}@{}}
\optlabel{A}\,P$\to$3, Q$\to$4, R$\to$2, S$\to$1 & \optlabel{B}\,P$\to$3, Q$\to$4, R$\to$1, S$\to$2 \\[3pt]
\optlabel{C}\,P$\to$4, Q$\to$3, R$\to$1, S$\to$2 & \optlabel{D}\,P$\to$3, Q$\to$1, R$\to$4, S$\to$2
\end{tabularx}
\end{qitem}
\qsep

\needspace{7\baselineskip}
\begin{qitem}{76}
Match each construction for an $n$-bit $N$ (List-I) with its logical-qubit count (List-II).
\par\smallskip\noindent{\renewcommand{\arraystretch}{1.35}%
\begin{tabularx}{\linewidth}{@{}c>{\raggedright\arraybackslash}p{0.33\linewidth}@{\hspace{1.2em}}cL@{}}
\toprule
\multicolumn{2}{@{}l}{\textbf{List-I}} & \multicolumn{2}{l@{}}{\textbf{List-II}} \\ \midrule
\textbf{(P)} & \texttt{Shor.ipynb} ($t = 2n$ counting $+$ $n$ work) & \textbf{(1)} & $n + 1$ \\
\textbf{(Q)} & Beauregard (2003) & \textbf{(2)} & $3n + 0.002\,n\log_2 n$ \\
\textbf{(R)} & Gidney--Ekerå (2019) & \textbf{(3)} & $3n$ \\
\textbf{(S)} & one recycled counting qubit $+$ $n$ work qubits with exact (oracle) multipliers & \textbf{(4)} & $2n + 3$ \\
\bottomrule
\end{tabularx}}
\par\smallskip\noindent\begin{tabularx}{\linewidth}{@{}*{2}{L}@{}}
\optlabel{A}\,P$\to$3, Q$\to$4, R$\to$2, S$\to$1 & \optlabel{B}\,P$\to$4, Q$\to$3, R$\to$2, S$\to$1 \\[3pt]
\optlabel{C}\,P$\to$3, Q$\to$4, R$\to$1, S$\to$2 & \optlabel{D}\,P$\to$3, Q$\to$2, R$\to$4, S$\to$1
\end{tabularx}
\end{qitem}
\qsep

\needspace{7\baselineskip}
\begin{qitem}{77}
\textbf{Assertion (A):} Migrating to post-quantum cryptography is urgent even though no current quantum computer can run Shor's algorithm at
cryptographic sizes.\\[1mm]
\textbf{Reason (R):} Shor's algorithm also solves the discrete-logarithm problem.
\begin{itemize}[leftmargin=2.6em, labelsep=0.4em, itemsep=2pt, topsep=3pt]
\item[\optlabel{A}] Both Assertion (A) and Reason (R) are true, and R is the correct explanation of A
\item[\optlabel{B}] Both A and R are true, but R is \emph{not} the correct explanation of A
\item[\optlabel{C}] A is true, but R is false
\item[\optlabel{D}] A is false, but R is true
\end{itemize}
\end{qitem}
\qsep

\needspace{7\baselineskip}
\begin{qitem}{78}
\textbf{Statement I:} In the notebook's construction, a $32$-bit $N$ needs $96$ qubits.\\[1mm]
\textbf{Statement II:} Repeated squaring lets $a^x \bmod N$ be implemented with $t$ controlled multiplications by the classically precomputed constants $a^{2^j} \bmod N$.
\begin{itemize}[leftmargin=2.6em, labelsep=0.4em, itemsep=2pt, topsep=3pt]
\item[\optlabel{A}] Both Statement I and Statement II are true
\item[\optlabel{B}] Both Statement I and Statement II are false
\item[\optlabel{C}] Statement I is true, but Statement II is false
\item[\optlabel{D}] Statement I is false, but Statement II is true
\end{itemize}
\end{qitem}
\qsep

\needspace{20\baselineskip}
\begin{passagebox}{Paragraph for Questions 79 and 80}
A gate-level version of the order-finding circuit for $N = 15$, $a = 7$ uses $t = 3$ counting qubits and $4$ work qubits. The
multiplications by $7$ and $4 = 7^2 \bmod 15$ are compiled into controlled SWAPs (each CNOT--Toffoli--CNOT, Toffoli $=$ 6 CNOTs) and CNOTs; the
inverse QFT uses $3$ CNOTs (one SWAP) and $3$ controlled phases. In total the circuit has $127$ CNOTs. It was simulated with the CUDA-Q
\texttt{density-matrix-cpu} target and a depolarizing noise model: a one-qubit channel with parameter $p$ after every $H$, $R_z$ and $X$, and a
two-qubit channel with parameter $p$ after every CNOT and controlled phase. The probability that one shot yields the factors $3 \times 5$ was
(re-executed, exact from the density matrix):
\begin{center}\begin{tabular}{lccccc}\toprule
$p$ & $0$ & $0.001$ & $0.003$ & $0.01$ & $0.03$\\\midrule
success probability & $0.750$ & $0.740$ & $0.726$ & $0.703$ & $0.665$\\\bottomrule
\end{tabular}\end{center}
If the counting register were completely random (uniform $y$), the success probability would be $0.375$. These are simulator results, not
measurements on a quantum processor.
\end{passagebox}
\needspace{7\baselineskip}
\begin{qitem}{79}
At $p = 0.01$, the probability that \emph{all} $127$ CNOTs are free of error is about $0.99^{127}$. Compare it with the success probability.
\begin{itemize}[leftmargin=2.6em, labelsep=0.4em, itemsep=2pt, topsep=3pt]
\item[\optlabel{A}] $0.99^{127} \approx 0.28$, yet the success probability is $0.703$: many errors do not change $y$ or still lead to a correct order
\item[\optlabel{B}] $0.99^{127} \approx 0.28$, so the success probability must be at most $0.28$
\item[\optlabel{C}] $0.99^{127} \approx 0.88$, equal to the success probability
\item[\optlabel{D}] $0.99^{127} \approx 0.01$, so the result is pure noise
\end{itemize}
\end{qitem}
\qsep

\needspace{7\baselineskip}
\begin{qitem}{80}
Which conclusion is supported by the table?
\begin{itemize}[leftmargin=2.6em, labelsep=0.4em, itemsep=2pt, topsep=3pt]
\item[\optlabel{A}] noise increases the success probability
\item[\optlabel{B}] the results show that a quantum processor factored $15$
\item[\optlabel{C}] the success probability decreases monotonically from $0.750$ towards the random baseline $0.375$; at $p = 0.03$ it is $0.665$, still above the baseline
\item[\optlabel{D}] at $p = 0.03$ the circuit is completely random
\end{itemize}
\end{qitem}
\qsep

\newpage

\begin{center}{\Large\bfseries\color{brand} RESPONSE SHEET}\\[2pt]
{\small Darken exactly ONE bubble per question. Leave all four bubbles empty to skip a question (0 marks).}\end{center}
\textbf{Name:} \hrulefill\quad \textbf{Roll No.:} \rule{35mm}{0.4pt}\par\medskip
\noindent\begin{minipage}[t]{0.235\linewidth}\centering
\textbf{Section A}\\[2pt]
\begin{tabular}{@{}r@{\hspace{5pt}}l@{}}
\textbf{1} & \bub{A}\,\bub{B}\,\bub{C}\,\bub{D} \\[1.6pt]
\textbf{2} & \bub{A}\,\bub{B}\,\bub{C}\,\bub{D} \\[1.6pt]
\textbf{3} & \bub{A}\,\bub{B}\,\bub{C}\,\bub{D} \\[1.6pt]
\textbf{4} & \bub{A}\,\bub{B}\,\bub{C}\,\bub{D} \\[1.6pt]
\textbf{5} & \bub{A}\,\bub{B}\,\bub{C}\,\bub{D} \\[1.6pt]
\textbf{6} & \bub{A}\,\bub{B}\,\bub{C}\,\bub{D} \\[1.6pt]
\textbf{7} & \bub{A}\,\bub{B}\,\bub{C}\,\bub{D} \\[1.6pt]
\textbf{8} & \bub{A}\,\bub{B}\,\bub{C}\,\bub{D} \\[1.6pt]
\textbf{9} & \bub{A}\,\bub{B}\,\bub{C}\,\bub{D} \\[1.6pt]
\textbf{10} & \bub{A}\,\bub{B}\,\bub{C}\,\bub{D} \\[1.6pt]
\textbf{11} & \bub{A}\,\bub{B}\,\bub{C}\,\bub{D} \\[1.6pt]
\textbf{12} & \bub{A}\,\bub{B}\,\bub{C}\,\bub{D} \\[1.6pt]
\textbf{13} & \bub{A}\,\bub{B}\,\bub{C}\,\bub{D} \\[1.6pt]
\textbf{14} & \bub{A}\,\bub{B}\,\bub{C}\,\bub{D} \\[1.6pt]
\textbf{15} & \bub{A}\,\bub{B}\,\bub{C}\,\bub{D} \\[1.6pt]
\textbf{16} & \bub{A}\,\bub{B}\,\bub{C}\,\bub{D} \\[1.6pt]
\textbf{17} & \bub{A}\,\bub{B}\,\bub{C}\,\bub{D} \\[1.6pt]
\textbf{18} & \bub{A}\,\bub{B}\,\bub{C}\,\bub{D} \\[1.6pt]
\textbf{19} & \bub{A}\,\bub{B}\,\bub{C}\,\bub{D} \\[1.6pt]
\textbf{20} & \bub{A}\,\bub{B}\,\bub{C}\,\bub{D} \\[1.6pt]
\end{tabular}
\end{minipage}\hfill\begin{minipage}[t]{0.235\linewidth}\centering
\textbf{Section B}\\[2pt]
\begin{tabular}{@{}r@{\hspace{5pt}}l@{}}
\textbf{21} & \bub{A}\,\bub{B}\,\bub{C}\,\bub{D} \\[1.6pt]
\textbf{22} & \bub{A}\,\bub{B}\,\bub{C}\,\bub{D} \\[1.6pt]
\textbf{23} & \bub{A}\,\bub{B}\,\bub{C}\,\bub{D} \\[1.6pt]
\textbf{24} & \bub{A}\,\bub{B}\,\bub{C}\,\bub{D} \\[1.6pt]
\textbf{25} & \bub{A}\,\bub{B}\,\bub{C}\,\bub{D} \\[1.6pt]
\textbf{26} & \bub{A}\,\bub{B}\,\bub{C}\,\bub{D} \\[1.6pt]
\textbf{27} & \bub{A}\,\bub{B}\,\bub{C}\,\bub{D} \\[1.6pt]
\textbf{28} & \bub{A}\,\bub{B}\,\bub{C}\,\bub{D} \\[1.6pt]
\textbf{29} & \bub{A}\,\bub{B}\,\bub{C}\,\bub{D} \\[1.6pt]
\textbf{30} & \bub{A}\,\bub{B}\,\bub{C}\,\bub{D} \\[1.6pt]
\textbf{31} & \bub{A}\,\bub{B}\,\bub{C}\,\bub{D} \\[1.6pt]
\textbf{32} & \bub{A}\,\bub{B}\,\bub{C}\,\bub{D} \\[1.6pt]
\textbf{33} & \bub{A}\,\bub{B}\,\bub{C}\,\bub{D} \\[1.6pt]
\textbf{34} & \bub{A}\,\bub{B}\,\bub{C}\,\bub{D} \\[1.6pt]
\textbf{35} & \bub{A}\,\bub{B}\,\bub{C}\,\bub{D} \\[1.6pt]
\textbf{36} & \bub{A}\,\bub{B}\,\bub{C}\,\bub{D} \\[1.6pt]
\textbf{37} & \bub{A}\,\bub{B}\,\bub{C}\,\bub{D} \\[1.6pt]
\textbf{38} & \bub{A}\,\bub{B}\,\bub{C}\,\bub{D} \\[1.6pt]
\textbf{39} & \bub{A}\,\bub{B}\,\bub{C}\,\bub{D} \\[1.6pt]
\textbf{40} & \bub{A}\,\bub{B}\,\bub{C}\,\bub{D} \\[1.6pt]
\end{tabular}
\end{minipage}\hfill\begin{minipage}[t]{0.235\linewidth}\centering
\textbf{Section C}\\[2pt]
\begin{tabular}{@{}r@{\hspace{5pt}}l@{}}
\textbf{41} & \bub{A}\,\bub{B}\,\bub{C}\,\bub{D} \\[1.6pt]
\textbf{42} & \bub{A}\,\bub{B}\,\bub{C}\,\bub{D} \\[1.6pt]
\textbf{43} & \bub{A}\,\bub{B}\,\bub{C}\,\bub{D} \\[1.6pt]
\textbf{44} & \bub{A}\,\bub{B}\,\bub{C}\,\bub{D} \\[1.6pt]
\textbf{45} & \bub{A}\,\bub{B}\,\bub{C}\,\bub{D} \\[1.6pt]
\textbf{46} & \bub{A}\,\bub{B}\,\bub{C}\,\bub{D} \\[1.6pt]
\textbf{47} & \bub{A}\,\bub{B}\,\bub{C}\,\bub{D} \\[1.6pt]
\textbf{48} & \bub{A}\,\bub{B}\,\bub{C}\,\bub{D} \\[1.6pt]
\textbf{49} & \bub{A}\,\bub{B}\,\bub{C}\,\bub{D} \\[1.6pt]
\textbf{50} & \bub{A}\,\bub{B}\,\bub{C}\,\bub{D} \\[1.6pt]
\textbf{51} & \bub{A}\,\bub{B}\,\bub{C}\,\bub{D} \\[1.6pt]
\textbf{52} & \bub{A}\,\bub{B}\,\bub{C}\,\bub{D} \\[1.6pt]
\textbf{53} & \bub{A}\,\bub{B}\,\bub{C}\,\bub{D} \\[1.6pt]
\textbf{54} & \bub{A}\,\bub{B}\,\bub{C}\,\bub{D} \\[1.6pt]
\textbf{55} & \bub{A}\,\bub{B}\,\bub{C}\,\bub{D} \\[1.6pt]
\textbf{56} & \bub{A}\,\bub{B}\,\bub{C}\,\bub{D} \\[1.6pt]
\textbf{57} & \bub{A}\,\bub{B}\,\bub{C}\,\bub{D} \\[1.6pt]
\textbf{58} & \bub{A}\,\bub{B}\,\bub{C}\,\bub{D} \\[1.6pt]
\textbf{59} & \bub{A}\,\bub{B}\,\bub{C}\,\bub{D} \\[1.6pt]
\textbf{60} & \bub{A}\,\bub{B}\,\bub{C}\,\bub{D} \\[1.6pt]
\end{tabular}
\end{minipage}\hfill\begin{minipage}[t]{0.235\linewidth}\centering
\textbf{Section D}\\[2pt]
\begin{tabular}{@{}r@{\hspace{5pt}}l@{}}
\textbf{61} & \bub{A}\,\bub{B}\,\bub{C}\,\bub{D} \\[1.6pt]
\textbf{62} & \bub{A}\,\bub{B}\,\bub{C}\,\bub{D} \\[1.6pt]
\textbf{63} & \bub{A}\,\bub{B}\,\bub{C}\,\bub{D} \\[1.6pt]
\textbf{64} & \bub{A}\,\bub{B}\,\bub{C}\,\bub{D} \\[1.6pt]
\textbf{65} & \bub{A}\,\bub{B}\,\bub{C}\,\bub{D} \\[1.6pt]
\textbf{66} & \bub{A}\,\bub{B}\,\bub{C}\,\bub{D} \\[1.6pt]
\textbf{67} & \bub{A}\,\bub{B}\,\bub{C}\,\bub{D} \\[1.6pt]
\textbf{68} & \bub{A}\,\bub{B}\,\bub{C}\,\bub{D} \\[1.6pt]
\textbf{69} & \bub{A}\,\bub{B}\,\bub{C}\,\bub{D} \\[1.6pt]
\textbf{70} & \bub{A}\,\bub{B}\,\bub{C}\,\bub{D} \\[1.6pt]
\textbf{71} & \bub{A}\,\bub{B}\,\bub{C}\,\bub{D} \\[1.6pt]
\textbf{72} & \bub{A}\,\bub{B}\,\bub{C}\,\bub{D} \\[1.6pt]
\textbf{73} & \bub{A}\,\bub{B}\,\bub{C}\,\bub{D} \\[1.6pt]
\textbf{74} & \bub{A}\,\bub{B}\,\bub{C}\,\bub{D} \\[1.6pt]
\textbf{75} & \bub{A}\,\bub{B}\,\bub{C}\,\bub{D} \\[1.6pt]
\textbf{76} & \bub{A}\,\bub{B}\,\bub{C}\,\bub{D} \\[1.6pt]
\textbf{77} & \bub{A}\,\bub{B}\,\bub{C}\,\bub{D} \\[1.6pt]
\textbf{78} & \bub{A}\,\bub{B}\,\bub{C}\,\bub{D} \\[1.6pt]
\textbf{79} & \bub{A}\,\bub{B}\,\bub{C}\,\bub{D} \\[1.6pt]
\textbf{80} & \bub{A}\,\bub{B}\,\bub{C}\,\bub{D} \\[1.6pt]
\end{tabular}
\end{minipage}
\par\bigskip
\begin{center}\textbf{Score computation} (filled in by the evaluator, using the Answer Book)\\[4pt]
\renewcommand{\arraystretch}{1.5}
\begin{tabular}{|c|>{\raggedright\arraybackslash}p{4.6cm}|c|c|c|c|}\hline
\textbf{Sec.} & \textbf{Topic} & \textbf{Attempted} & \textbf{Correct} $N_{\mathrm c}$ & \textbf{Wrong} $N_{\mathrm w}$ & \textbf{Marks} $3N_{\mathrm c} - N_{\mathrm w}$ \\ \hline
A & Mathematical and Number Theory Foundations & & & & \\ \hline
B & Quantum Computing Foundations & & & & \\ \hline
C & Shor's Algorithm Fundamentals & & & & \\ \hline
D & Advanced Shor's Algorithm and Implementation & & & & \\ \hline
\multicolumn{5}{|r|}{\textbf{Total (out of 240)}} & \\ \hline
\end{tabular}\end{center}

\end{document}